%-------------------------
% Resume in Latex
% Author : Jake Gutierrez
% Based off of: https://github.com/sb2nov/resume
% License : MIT
%------------------------

\documentclass[letterpaper,11pt]{article}

\usepackage{latexsym}
\usepackage[empty]{fullpage}
\usepackage{titlesec}
\usepackage{marvosym}
\usepackage[usenames,dvipsnames]{color}
\usepackage{verbatim}
\usepackage{enumitem}
\usepackage[hidelinks]{hyperref}
\usepackage{fancyhdr}
\usepackage[english]{babel}
\usepackage{tabularx}
\input{glyphtounicode}


%----------FONT OPTIONS----------
% sans-serif
% \usepackage[sfdefault]{FiraSans}
% \usepackage[sfdefault]{roboto}
% \usepackage[sfdefault]{noto-sans}
% \usepackage[default]{sourcesanspro}

% serif
% \usepackage{CormorantGaramond}
% \usepackage{charter}


\pagestyle{fancy}
\fancyhf{} % clear all header and footer fields
\fancyfoot{}
\renewcommand{\headrulewidth}{0pt}
\renewcommand{\footrulewidth}{0pt}

% Adjust margins
\addtolength{\oddsidemargin}{-0.5in}
\addtolength{\evensidemargin}{-0.5in}
\addtolength{\textwidth}{1in}
\addtolength{\topmargin}{-.5in}
\addtolength{\textheight}{1.0in}

\urlstyle{same}

\raggedbottom
\raggedright
\setlength{\tabcolsep}{0in}

% Sections formatting
\titleformat{\section}{
  \vspace{-4pt}\scshape\raggedright\large
}{}{0em}{}[\color{black}\titlerule \vspace{-5pt}]

% Ensure that generate pdf is machine readable/ATS parsable
\pdfgentounicode=1

%-------------------------
% Custom commands
\newcommand{\resumeItem}[1]{
  \item\small{
    {#1 \vspace{-2pt}}
  }
}

\newcommand{\resumeSubheading}[4]{
  \vspace{-2pt}\item
    \begin{tabular*}{0.97\textwidth}[t]{l@{\extracolsep{\fill}}r}
      \textbf{#1} & #2 \\
      \textit{\small#3} & \textit{\small #4} \\
    \end{tabular*}\vspace{-7pt}
}

\newcommand{\resumeSubSubheading}[2]{
    \item
    \begin{tabular*}{0.97\textwidth}{l@{\extracolsep{\fill}}r}
      \textit{\small#1} & \textit{\small #2} \\
    \end{tabular*}\vspace{-7pt}
}

\newcommand{\resumeProjectHeading}[2]{
    \item
    \begin{tabular*}{0.97\textwidth}{l@{\extracolsep{\fill}}r}
      \small#1 & #2 \\
    \end{tabular*}\vspace{-7pt}
}

\newcommand{\resumeSubItem}[1]{\resumeItem{#1}\vspace{-4pt}}

\renewcommand\labelitemii{$\vcenter{\hbox{\tiny$\bullet$}}$}

\newcommand{\resumeSubHeadingListStart}{\begin{itemize}[leftmargin=0.15in, label={}]}
\newcommand{\resumeSubHeadingListEnd}{\end{itemize}}
\newcommand{\resumeItemListStart}{\begin{itemize}}
\newcommand{\resumeItemListEnd}{\end{itemize}\vspace{-5pt}}

%-------------------------------------------
%%%%%%  RESUME STARTS HERE  %%%%%%%%%%%%%%%%%%%%%%%%%%%%


\begin{document}

%----------HEADING----------
% \begin{tabular*}{\textwidth}{l@{\extracolsep{\fill}}r}
%   \textbf{\href{http://sourabhbajaj.com/}{\Large Sourabh Bajaj}} & Email : \href{mailto:sourabh@sourabhbajaj.com}{sourabh@sourabhbajaj.com}\\
%   \href{http://sourabhbajaj.com/}{http://www.sourabhbajaj.com} & Mobile : +1-123-456-7890 \\
% \end{tabular*}

\begin{center}
    \textbf{\Huge \scshape James Cheng} \\ \vspace{1pt}
    \small +1 (408)-981-8876 $|$ \href{mailto:jzcheng@stanford.edu}{\underline{jzcheng@stanford.edu}} $|$
    \href{https://www.linkedin.com/in/jamescheng24/}{\underline{linkedin.com/in/jamescheng24/}}
\end{center}


%-----------EDUCATION-----------
\section{Education}
  \resumeSubHeadingListStart
    \resumeSubheading
      {Stanford University}{Stanford, CA}
      {M.S. in Computer Science, Artificial Intelligence Specialization}{Expected Jun. 2026}
      \resumeItemListStart
        \resumeItem{\textbf{Courses}: Reinforcement Learning, NLP, Computer Vision, Robot Autonomy \& Perception, Parallel Computing}
        \resumeItemListEnd
    \resumeSubheading
      {University of California, Berkeley - Regents' and Chancellor's Scholar}{Berkeley, CA}
      {B.A. in Computer Science and B.S. in Business Administration, GPA: 3.92/4.00}{Aug. 2020 -- May 2024}
      \resumeItemListStart
        \resumeItem{\textbf{Courses}: Algorithms, Operating Systems, Machine Learning, Data Structures, Random Processes, Optimization}
      \resumeItemListEnd
  \resumeSubHeadingListEnd


%-----------EXPERIENCE-----------
\section{Experience}
  \resumeSubHeadingListStart
    \resumeSubheading
        {Stanford CS}{Stanford, CA}
        {Research Engineer - The Movement Lab}{Sep. 2025 --- Present}
        \resumeItemListStart
            \resumeItem{Porting advanced motion imitation methods (DeepMimic, AMP, ASE, ADD) from MimicKit to MuJoCo Lab, integrating them with the Isaac Lab API for scalable humanoid motion control and policy learning}
        \resumeItemListEnd

        \resumeSubSubheading
        {Teaching Assistant (CS234: Reinforcement Learning, CS161: Design \& Analysis of Algorithms)}{Sep. 2025 -- Present}
        \resumeItemListStart
            \resumeItem{Selected as Course Staff for Graduate Reinforcement Learning (CS 234), guiding students through implementations of PPO, RLHF, and DPO in PyTorch/MuJoCo.}
            \resumeItem{Led office hours for Algorithms, clarifying complex topics in dynamic programming, graph theory, and randomized algorithms for 300+ students.}
        \resumeItemListEnd

    \resumeSubheading
      {Amazon}{Sunnyvale, CA}
      {Software Engineering Intern - Kindle}{Jun. 2025 -- Sep. 2025}
      \resumeItemListStart
        \resumeItem{Architected the backend integration for on-device AI image generation, orchestrating low-latency communication between constrained Kindle hardware (C++) and cloud inference (AWS Bedrock/Lambda)}
      \resumeItemListEnd

    \resumeSubheading
      {UC Berkeley EECS}{Berkeley, CA}
      {Undergraduate Research Assistant - BAIR Lab (Advised by Prof. Shankar Sastry)}{Jun. 2023 -- May 2024}
      \resumeItemListStart
        \resumeItem{Implemented follower-agnostic convergent algorithm for a traffic routing Stackelberg Game using Python; ran numerical experiments across 12 trajectories and tuned step-size, perturbation, and regularization parameters, and visualized convergence of objective values and tolls to stationary point; \textbf{paper accepted to CDC 2024}}
        % \resumeItem{Developed program to simulate arc-based traffic assignment model using Python and visualized simulation}
      \resumeItemListEnd

      \resumeSubSubheading
        {Teaching Assistant (CS 70: Discrete Mathematics and Probability)}{Aug. 2022 -- May 2024}
      \resumeItemListStart
        \resumeItem{Led students in discussion, assisted in office hours, and graded homework for 500+ students on topics such as proofs, stable matching, graph theory, modular arithmetic, cryptography, computability, combinatorics, and probability}
      \resumeItemListEnd


% -----------Multiple Positions Heading-----------
%    \resumeSubSubheading
%     {Software Engineer I}{Oct 2014 - Sep 2016}
%     \resumeItemListStart
%        \resumeItem{Apache Beam}
%          {Apache Beam is a unified model for defining both batch and streaming data-parallel processing pipelines}
%     \resumeItemListEnd
%    \resumeSubHeadingListEnd
%-------------------------------------------

    \resumeSubheading
      {RUCKUS Networks - CommScope}{Sunnyvale, CA}
      {Software Engineering Intern - Application Development Engineering Team}{Jun. 2023 -- Aug. 2023}
      \resumeItemListStart
          \resumeItem{Developed Airbyte connectors in Python to ingest 10k+ daily records into BigQuery and built CI/CD pipelines with Docker/Bitbucket, reducing manual deployment overhead.}
        % \resumeItem{Developed 2 Airbyte data connectors using Python for third-party services by calling REST APIs, resulting in 10k+ records ingested into BigQuery daily and \textbf{reducing time to pull 1M records from 5 hours to a few seconds}}
        % \resumeItem{Implemented 30+ Pytest unit tests and created Dockerfiles with Bitbucket CI/CD pipelines to automate testing and deployment, reducing manual overhead.}
        % \resumeItem{Implemented 30+ unit tests in Pytest to ensure proper functionality of each file and function and wrote SQL and YAML for data transformations to integrate dbt with Airbyte}
        % \resumeItem{Created Dockerfile to build Docker image for integration with Airbyte along with developing CI/CD build with Bitbucket pipelines for automatically building and testing application during each code push}
        % \resumeItem{Analyzed ingested data using PowerBI by generating 10 visualizations for activity metrics}
    \resumeItemListEnd

    % \resumeSubheading
    %   {Paramount Global}{New York, NY}
    %   {Data Engineering Intern - Data Technology Solutions Team}{Jun. 2022 -- Aug. 2022}
    %   \resumeItemListStart
    %     \resumeItem{Implemented 6 types of data quality tests in SQL and YAML via dbt (Data Build Tool) for 20+ models: freshness, null check, uniqueness, referential integrity, consistency, and accuracy}
    %     \resumeItem{Automated and orchestrated data testing through writing DAG files in Python for Apache Airflow}
    %   \resumeItemListEnd

      % \resumeSubheading
      % {Align Technology}{San Jose, CA}
      % {Data Analytics Intern - Global Business Analytics Team}{Jun. 2021 -- Aug. 2021}
      % \resumeItemListStart
      %   \resumeItem{Analyzed global data using SQL Server Management Studio and Microsoft Excel, gaining insights into financial KPIs such as revenue, market penetration, and profit margin trends across customer demographics}
      %   \resumeItem{Researched and discovered \$17 M incremental revenue per year opportunity in Invisalign Doctor Service WebStore}

      % \resumeItemListEnd

  \resumeSubHeadingListEnd


%-----------PROJECTS-----------
\section{Projects}
    \resumeSubHeadingListStart
        % \resumeProjectHeading
        %     {\textbf{Automatic Music Transcription: Sustain Pedal Prediction} $|$ \emph{Python, PyTorch}}{\textit{Sep. 2024 -- Dec. 2024}}
        %     \resumeItemListStart
        %         \resumeItem{Contributed to CNN and LSTM model that takes as input audio piano music and outputs predictions of whether the sustain pedal is used for each time frame, achieving \textbf{92\% accuracy}}
        %         \resumeItem{Developed evaluation pipeline and logistic regression baseline with cross-validation which achieved 68\% accuracy}
        %     \resumeItemListEnd

      %   \resumeProjectHeading
      %     {\textbf{Stock Prediction} $|$ \emph{Python, Scikit-learn, Pandas, NumPy}}{\textit{Jan. 2024 -- May 2024}}
      %     \resumeItemListStart
      %       \resumeItem{Transformed 2017-2018 company data via imputation and PCA to predict growth and create profitable portfolio}
      %       \resumeItem{Cross-validated across linear regression, kNN, neural network, decision tree, and naive Bayes models to develop a \textbf{top-3 performing model in terms of profit out of 52 students}}
      % \resumeItemListEnd
      \resumeProjectHeading
  {\textbf{Constrained Text Generation Benchmark (CTG)} $|$ \emph{Python, uv, Together API, Anthropic API, spaCy}}{\textit{Jan. 2026 -- Present}}
  \resumeItemListStart
    \resumeItem{Designed a \textbf{348-prompt} benchmark (101 local / 97 global / 150 hybrid constraints) with dual-path verification (programmatic + LLM grading) to quantify autoregressive models' constraint-following limitations across 6 LLMs.}
    \resumeItem{Evaluated Claude, Llama, and DeepSeek models on \textbf{7,889 expanded prompt variants}, confirming a steep difficulty gradient (local 37--53\% vs. global 12--21\% pass rates) driven by autoregressive planning bottlenecks.}
    \resumeItem{Built SFT and DPO fine-tuning pipelines on Together AI to improve constraint adherence, plus meta-ablation analyses comparing hybrid composition pass rates against individual constraint baselines.}
  \resumeItemListEnd


      \resumeProjectHeading
          {\textbf{LLM-as-Planner vs RL on Procgen} $|$ \emph{PyTorch, StableBaselines3, Amazon Bedrock}}{\textit{Mar. 2025 -- Jun. 2025}}
          \resumeItemListStart
            \resumeItem{Built hybrid RL-LLM agents with entropy-based mode switching, boosting task success rates up to 2× by enabling sustained LLM-guided control in complex environments.}
            \resumeItem{Benchmarked LLM planning vs. RL baselines (PPO, DQN) across procedurally generated environments, uncovering trade-offs between zero-shot generalization and fine-grained control.}
          \resumeItemListEnd

      % \resumeProjectHeading
      %     {\textbf{RUCKUS GPT Chatbot} $|$ \emph{Python, LangChain, OpenAI, Streamlit, FAISS, PyPDF}}{\textit{Jun. 2023 -- Aug. 2023}}
      %     \resumeItemListStart
      %       \resumeItem{Developed customer support chatbot in Python with embedded documents and memory to have complex conversations, achieving \textbf{top 5 out of 300+ projects globally in hackathon}}
      %       \resumeItem{Integrated the following Python packages: Streamlit for user interface, OpenAI for language model and embeddings, LangChain for conversation chain and chat history, FAISS for vector indexing, and PyPDF for parsing PDFs}
      %     \resumeItemListEnd
    \resumeSubHeadingListEnd



%
%-----------PROGRAMMING SKILLS-----------
\section{Technical Skills}
 \begin{itemize}[leftmargin=0.15in, label={}]
    \small{\item{
     \textbf{Languages/Packages}{: Python, C/C++ and CUDA, Rust, Java, SQL, PyTorch, NumPy, Pandas} \\
     \textbf{Platform/Tools}{: MuJoCo, VS Code, AWS, Git, Docker, Apache Airflow, Google Cloud Platform, Jupyter, PowerBI} \\
    }}
 \end{itemize}


%-------------------------------------------
\end{document}
